Un servicio REST no guarda estado entre peticiones \cite{fielding2000}, así que lo que decide entre consistencia y disponibilidad es la base de datos que tiene detrás. Con PACELC se ven dos momentos: si hay partición se elige A o C, y si no la hay se elige entre latencia y consistencia \cite{abadi2012}.

DynamoDB lo deja en manos del cliente: por defecto lee de forma eventual, que cuesta la mitad, y si el cliente manda \texttt{ConsistentRead=true} recibe el dato más reciente \cite{aws_consistencia}. La caché de HTTP hace algo parecido: entrega la copia guardada mientras siga fresca y con un \texttt{ETag} pregunta si cambió \cite{rfc9111}. Nuestro PFC tiene una sola base PostgreSQL, así que siempre lee lo último, pero si la base no responde el servicio falla: prefiere C antes que A. Abadi deja claro que no se escoge \enquote{dos de tres} para siempre, porque las particiones son raras \cite{abadi2012}.
